\chapter{All finite symmetric polynomial nulls conserve content}\label{ch:pcall}
\section{Removing the homogeneity restriction}
The earlier content proof used $Q(gv)=g^dQ(v)$. A finite nonhomogeneous polynomial does not have one scaling degree. In principle, cancellation between its homogeneous pieces might make it universally null while allowing individual pieces to have nonzero sums. The following theorem excludes that escape.

\begin{theorem}\label{thm:pcall}
Let $Q$ be any finite symmetric polynomial over $\C$, restricted to $\sum_i v_i=0$. Suppose
\[
\sum_{v\in\Vn}W(v)Q(v)=0
\]
for every positive integer $n$. Write $Q=\sum_dQ_d$ in homogeneous degrees. Then every $Q_d$ is separately universally null. Moreover, the weighted $Q$-sum is zero in each exact-content sector of every shell.
\end{theorem}

The finite-polynomial hypothesis is substantive. The argument below separates finitely many modular weights. It says nothing about an arbitrary infinite series, arbitrary functions on the packets or nonsymmetric corrections. For a sign consequence, corrected weights must be real, as they are for the retained rational corrections. Complex coefficients have no ordered sign.

\section{Harmonic decomposition of an alternating polynomial}
For a homogeneous $Q_d$, put $P=DQ_d$ and $M=10+d$. The standard decomposition into harmonic components on the rank-four Euclidean hyperplane is
\begin{equation}\label{eq:harmdecomp}
P=\sum_{j\ge0}s_2^jh_{M-2j},
\end{equation}
where $h_\ell$ is harmonic homogeneous of degree $\ell$. It can be constructed recursively from the Laplacian. For a harmonic $h_\ell$,
\begin{equation}\label{eq:radiallap}
\Delta_H(s_2^jh_\ell)=4j(\ell+j+1)s_2^{j-1}h_\ell.
\end{equation}
The nonzero coefficients in this formula give uniqueness of the decomposition. The operator commutes with coordinate permutations, so every $h_\ell$ is alternating when $P$ is alternating. Hence $D$ divides every nonzero $h_\ell$, and $\ell\ge10$. Thus only $j\le d/2$ can occur.

When $d$ is odd, the shell sum vanishes directly under $J$ as in \eqref{eq:oddnull}, separately inside each content class. We may therefore concentrate on even $d$, when all harmonic degrees in \eqref{eq:harmdecomp} are even and the scalar harmonic theta argument applies.

\section{Radial factors become derivatives}
For an alternating harmonic $h_\ell$, write
\[
\theta_h(\tau)=\sum_{\substack{v_i\equiv i\ (5)\\\sum v_i=0}}
h(v)\exp\left(2\pi i\tau\frac{\|v\|^2}{10}\right).
\]
Chapter~\ref{ch:theta} establishes that it is a scalar cusp form of weight $\ell+2$. The normalization includes only a fixed scalar if one divides by $288$.

Each application of $\DQ=q\,d/dq$ multiplies the shell coefficient by $\|v\|^2/10$. Therefore
\begin{equation}\label{eq:radialderivative}
\theta_{s_2^jh}=10^j\DQ^j\theta_h.
\end{equation}
Differentiating a modular transformation law shows that $\DQ^jf$, for modular weight $w$, is quasimodular of weight $w+2j$ and depth at most $j$. Combining \eqref{eq:harmdecomp} and \eqref{eq:radialderivative}, the homogeneous shell series
\[
F_d(\tau)=\sum_{n\ge1}\left(\sum_{v\in\Vn}W(v)Q_d(v)\right)q^n
\]
is quasimodular of weight
\[
(M-2j+2)+2j=12+d,
\]
independent of the radial index $j$. Its depth is at most $d/2$. Polynomial growth and Gaussian decay justify the termwise differentiations.

The common weight for all terms within one homogeneous degree is the key. Different homogeneous degrees produce different weights. A nonhomogeneous null identity would have to make those different weights cancel, and the next step shows that they cannot.

\section{Separating the quasimodular weights}
We need only an elementary weight-separation statement. Suppose finitely many quasimodular functions $F_k$ of distinct nonnegative integer weights satisfy
\[
F_k(\gamma z)=\sum_{r=0}^{p_k}f_{k,r}(z)c^r(cz+d)^{k-r},
\qquad f_{k,0}=F_k,\qquad p_k\le k,
\]
for $\gamma=\left[\begin{smallmatrix}a&b\\c&d\end{smallmatrix}\right]\in SL_2(\Z)$. Constants from the usual quasimodular convention can be absorbed into the coefficient functions. Assume $\sum_kF_k=0$ identically.

Fix $z$ in the upper half-plane. The polynomial
\begin{equation}\label{eq:weightpolynomial}
P_z(c,d)=\sum_k\sum_r f_{k,r}(z)c^r(cz+d)^{k-r}
\end{equation}
vanishes at every relatively prime integer pair $(c,d)$, because every such pair is the bottom row of a modular matrix. These pairs are Zariski dense. To see this without a geometric theorem, fix any positive prime $c$. Infinitely many integers $d$ are coprime to it, so the polynomial in $d$ vanishes identically. Each of its coefficient polynomials in $c$ then vanishes at infinitely many primes, so is identically zero. Hence $P_z$ is the zero polynomial.

Every summand belonging to weight $k$ is homogeneous of total degree $k$ in $(c,d)$. Different homogeneous components of a zero polynomial vanish separately. Evaluate the degree-$k$ component at $(c,d)=(0,1)$. Only $r=0$ survives, giving $F_k(z)=0$. Since $z$ was arbitrary, every $F_k$ is zero.

Apply this to the finite even homogeneous degrees with weights $12+d$. Universal nullness of $Q$ says their sum is zero, after the odd pieces have already vanished. Thus every $F_d$ is zero. Equality of Fourier coefficients proves that every homogeneous $Q_d$ is universally null.

\section{Finishing the content proof}
Each homogeneous piece now satisfies Theorem~\ref{thm:homcontent}. Its exact-content contribution is zero on every shell. Summing the finitely many pieces proves the full nonhomogeneous claim.

The argument has two independent mechanisms. Modular transformation laws prevent different homogeneous degrees from canceling globally. The arithmetic scaling identity prevents a homogeneous null from moving mass between exact contents. Both are needed. Merely applying the homogeneous scaling formula to a nonhomogeneous $Q$ would skip the first mechanism.

\section{The strengthened composite boundary}
At $n=121$, the opposed content masses remain conserved under every finite symmetric polynomial universal-null correction. This is now true without a degree cutoff, a prescribed finite family or an initial homogeneity assumption. Any real corrected weights of one common sign would force both content totals to have that sign, contradicting the exact values.

The theorem closes one possible escape from the content invariant. It does not constrain all nonpolynomial operations, and it does not produce a new prime invariant. At a prime there is still one content class, and its total is still $\tau(p)$.

\chapter{Interpolation is not universal nullness}\label{ch:interpolation}
\section{Coordinates that identify exactly the orbits}
Fix a nonempty shell and put $C=10n$. Consider the normalized even symmetric invariants
\begin{equation}\label{eq:invariantmap}
I(v)=\left(\frac{s_4}{C^2},\frac{s_3^2}{C^3},
\frac{s_3s_5}{C^4},\frac{s_5^2}{C^5}\right).
\end{equation}
They are unchanged by $J$.

\begin{proposition}\label{prop:invariantseparation}
On a fixed shell, $I(v)=I(w)$ if and only if $w=v$ or $w=Jv$.
\end{proposition}
\begin{proof}
Equality gives equal $s_4$ and equal rank-one outer products of the real vectors $(s_3(v),s_5(v))$ and $(s_3(w),s_5(w))$. Those vectors agree up to a common sign, including the zero case. Thus the first five power sums of $v$ and either $w$ or $-w$ agree: $s_1=0$, $s_2=C$ and the remaining three by the preceding comparison.

Newton's identities recover the elementary symmetric polynomials from those power sums. The coordinate multisets therefore agree. Each multiset has exactly one coordinate in each residue class modulo $5$, so equal multisets imply $v=w$. Negated multisets imply $w_i=-v_{-i\bmod5}=J(v)_i$. The converse follows from the even invariance of all four coordinates in \eqref{eq:invariantmap}.
\end{proof}

This removes a potential geometric explanation for an evaluation defect. If an all-degree universal-null family has an additional annihilator on a shell, it cannot arise from two distinct $J$-orbits secretly having the same even symmetric coordinates. It must reflect a further arithmetic or modular restriction.

\section{An explicit homogeneous interpolation basis}
Let $v_1,\ldots,v_N$ represent the distinct $J$-orbits. For an integer $t\ge0$, define the degree-$10$ symmetric polynomial
\begin{equation}\label{eq:et}
E_t=s_4s_2^3+t\,s_3^2s_2^2+t^2s_3s_5s_2+t^3s_5^2.
\end{equation}
For two distinct orbits, $E_t(v_i)-E_t(v_j)$ is a nonzero polynomial in $t$ of degree at most $3$, by Proposition~\ref{prop:invariantseparation}. It has at most three roots. Among all pairs there are at most $3N(N-1)/2$ forbidden values. Hence some integer
\[
0\le t\le\frac{3N(N-1)}2
\]
separates all $e_i=E_t(v_i)$.

For that $t$, set
\begin{equation}\label{eq:li}
L_i(v)=\prod_{j\ne i}
\frac{E_t(v)-e_j(s_2(v)/C)^5}{e_i-e_j}.
\end{equation}
Each factor is homogeneous of degree $10$, and the denominators are nonzero rational numbers. The resulting polynomial is symmetric homogeneous of degree $10(N-1)$. On the shell, $s_2/C=1$, so
\[
L_i(v_j)=\delta_{ij}.
\]
Consequently, any prescribed rational orbit values $a_1,\ldots,a_N$ are interpolated by $Q=\sum_ia_iL_i$. If $N=1$, an ordinary constant suffices.

This is a constructive interpolation theorem, with a degree bound and explicit basis. It imposes neither harmonicity nor universal nullness.

\section{A shell-null interpolant can fail on the first shell}
Let $m_i$ be the nonzero mass of the $i$th orbit. Interpolation makes the map $Q\mapsto(m_iQ(v_i))$ onto $\Q^N$. Thus any zero-total transfer on this one shell can be produced by a homogeneous polynomial.

The condition that its weighted sum vanish on this shell is weaker than requiring it to vanish on every shell. The natural interpolation that kills all but the first orbit at $n=7$ has degree $20$ and zero correction total at that target. Its corrected survivor is $-16744$. Yet its weighted correction sum on shell $1$ is
\[
-\frac{449}{42371287350}\ne0.
\]
The analogous degree-$150$ construction at $31$ also fails already at shell $1$. These exact failures occur at small shells where different, universal nulls are known to supply sign certificates.

Thus ordinary interpolation cannot be substituted for the modular kernel construction. Both use polynomials, but they solve different systems of conditions.

\section{Harmonic projection does not preserve values}
One might try to interpolate first and then take a harmonic projection. There is no general principle preserving the target-shell values under this operation. Since $D$ is harmonic of degree $10$ on a rank-four space,
\[
\Delta_H(s_2D)=48D.
\]
The degree-$12$ top harmonic component of $s_2D$ is zero: the entire polynomial is a radial factor times a lower-degree harmonic polynomial. Yet $s_2D$ is nonzero at every packet on a nonempty shell.

More generally, the top harmonic projection of $s_2^kD$ at degree $10+2k$ is zero for every $k>0$. Radial information can disappear. Replacing an interpolant by its top harmonic piece may therefore destroy exactly the values it was designed to have.

Dimension growth does not repair this missing implication. Even if harmonic-polynomial dimensions grow quadratically while the number of modular coefficient constraints grows linearly, one must prove independence between the evaluation rows and the modular rows. A large kernel need not evaluate surjectively onto a prescribed finite subspace. The finite saturation profiles in Chapter~\ref{ch:transport} are useful examples, not that uniform independence theorem.

\section{The strongest conceivable content fullness}
Let the orbits be grouped by exact content and let $T_g$ be the mass of each group. Suppose, as an unproved hypothesis, that all-degree universal null transfers span every vector with zero sum in each content class. Then a strictly positive corrected mass on every orbit exists exactly when every nonempty $T_g$ is positive. Necessity is conservation. For sufficiency, distribute $T_g$ equally among its orbits, then use fullness to realize the difference from the original mass vector.

The negative version is identical. For a weakly nonnegative certificate with at least one strict mass, the criterion is $T_g\ge0$ for every class and $\sum_gT_g>0$.

At a prime there is one class. Under even this maximal fullness hypothesis, the existence of a common-sign certificate is therefore equivalent to $\tau(p)\ne0$. A concrete independently verified certificate can prove nonvanishing. The mere availability of every conserved redistribution cannot decide whether the conserved total vanishes.

The all-degree universal-null fullness theorem itself remains unproved in this program. Its proof would answer an interesting geometric and modular question. It would still leave an arithmetic sign or residue input to complete Lehmer.
